\documentclass[a4paper, 12pt]{article}
\usepackage{graphicx}
\usepackage{amsmath}
\usepackage{hyperref}
\usepackage{bookmark}
\usepackage{geometry}
\usepackage{amssymb}
\usepackage{helvet}
\usepackage{subcaption}

\renewcommand{\familydefault}{\sfdefault}
\geometry{left=0.75in, right=0.5in, top=1.0in, bottom=0.75in}
\title{Water Scarcity Hackathon - Phase 2 - Policy Optimization}
\author{Dami Akinniyi \textit{\(@yeleon\)}}
\date{\today}
\begin{document}
\maketitle
\section*{Abstract}
This report presents the design and evaluation of adaptive policies within a simulated 
multi-agent environment, focusing on ecofairness and ecological sustainability. 
The proposed framework integrates tunable quotas and incentive mechanisms to regulate 
agent behavior, promoting a dynamic balance of cooperation and competition while responding 
to environmental feedback. %Agents evolve under policy-guided rules, enabling exploration of 
%cooperative structures and system resilience under resource scarcity.
\\ \\
We examine how resource constraints and policy choices shape system dynamics. 
Prioritizing certain actors creates economic imbalances, heightens competition and increases ecological impact. 
While subsidies can promote cooperation and ecological alignment, excessive support weakens resilience in stressed environments. 
Adaptive, well-calibrated policies enhance robustness, ecological survival, and stability under uncertainty

\section{Introduction}
This project was developed as part of Phase 2 of the 2025 Streamflow Forecasting and Water 
Resource Regulation Hackathon. While Phase 1 focused on forecasting streamflow, Phase 2 
centers on the design and optimization of policies to govern the allocation of shared 
water resources in a game theory simulation environment.
\\ \\
Effective water policy must balance ecological protection, equitable access, and economic stability. 
Our approach emphasizes fairness—allocating water by priority—while ensuring sustainability for downstream ecosystems and resilience under uncertainty.
We iteratively test the simulation environment, manually adjusting key parameters (e.g., demand and penalty multipliers) 
to understand how policy rules affect economic and ecological outcomes. 
\\ \\
We observed the incentive function shapes drive cooperation and economic stability, while quota allocations and actor priorities impact ecological health.
Based on these insights, we design adaptive policies that respond to environmental signals, guiding agents toward 
sustainable behavior. Our framework includes:
\begin{itemize}
    \item Simulation Environment Parallelization
    \item Quota Function: A priority-aware allocation mechanism with configurable thresholds.
    \item Incentive Function: A fine/subsidy structure with both stepwise and smooth variants.
    \item Multiobjective Optimization : A Pareto-front approach using Optuna to discover trade-offs between economic efficiency and ecological viability.
    \item Policy Regulator Framework: A python class that encapsulates the policy design and optimization process to yield simulation-ready policy functions.  
\end{itemize}
\noindent This framework allows us to explore how adaptive, minimalistic rules can guide cooperation and robustness in a complex, uncertain environment.
\section{Simulation Environment Parallelization}
At the heart of the policy evaluation workflow is the $multi\_scenario\_analysis$ function, responsible 
for executing simulations across a range of behavioral and environmental configurations. This function 
plays a central role in assessing the robustness and generalizability of policy designs, ensuring that 
any proposed strategy is evaluated across diverse and uncertain scenarios. 
\\ \\
The original function is effective in exploring these scenarios sequentially. The parallelization of the 
$multi\_scenario\_analysis$ function is a crucial step in optimizing the policy evaluation process. 
By leveraging Python's concurrency capabilities, we can execute multiple scenarios simultaneously, 
significantly reducing the time required for each evaluation cycle. This is particularly important in the 
context of policy optimization, where rapid feedback on policy performance is essential for iterative 
refinement and decision-making.
\\ \\
We also enhance the function's flexibility to allow specifying custom scenario subsets—enabling 
targeted evaluation and more efficient use of compute resources during optimization. These enhancements 
integrate seamlessly with optuna’s multi-objective optimization framework, allowing for large-scale policy search.
\\ \\
Together, these improvements reduce the runtime of a 700 turns, 10 iterations simulation from approximately 
3 hours to less than 20 minutes per full evaluation cycle.
% \begin{itemize}
%     \item Simulation Environment Efficiency - Parallelizing functions to speed up implementation to enable full optimization.
%     \item Quota Function Design: Gathering historical streamflow data, water demand, and environmental factors.
%     \item Incentives Function Design: Using machine learning models to predict future streamflow based on historical data.
%     \item Policy Optimization: Developing the quota and incentives functions to allocate water resources effectively.
%     \item Evaluation: Assessing the performance of the policy using various metrics such as water scarcity index and user satisfaction.
%     \item Iteration: Refining the policy based on feedback and performance metrics.
% \end{itemize}
\section{Quota Function Design}
In water-stressed systems, resource allocation must carefully balance fairness, actor resilience, and 
ecological resilience — ensuring that critical needs are met without compromising the long-term health 
of the environment. This becomes especially important during crisis conditions, where water scarcity 
intensifies and allocation decisions carry systemic consequences.
\\ \\
\noindent To address this challenge, we propose a quota policy function for allocating water across multiple actors 
with varying demands and priorities. 

\subsection{Policy Mechanism}
The function is guided by three key parameters:

\begin{itemize}
    \item \textbf{rate ($r$)}: The baseline proportion of each actor's demand that is fulfilled during non-crisis periods. 
    This controls total water use and supports ecosustainability by preventing excessive extraction even under normal conditions.
    \item \textbf{priority\_effect ($p_e$)}: A tuning parameter that controls how much high-priority actors are favored during crises.
    When priority effect is 0, non-low-priority actors share water equally (relative to demand). Higher values increasingly concentrate 
    allocation toward the highest-priority actors. Low-priority actors receive no water during alert and crises, based on prior analysis 
    showing they maintain high buffer-to-demand ratios and can self-sustain through temporary shortages.
    \item \textbf{crisis\_rate ($r_c$)}: The overall reduction in water availability during crisis or alert conditions, defining what 
    percentage of each actor’s demand (before applying priority) is considered for allocation. This simulates external constraints and enforces 
    ecosustainability by capping water use according to environmental thresholds.
\end{itemize}  
\subsection{Mathematical Formulation}
Let:
\begin{itemize}
  \item $\mathbf{D} = (D_1, D_2, \dots, D_n)$: vector of average water pumped per actor.
  \item $\mathbf{P} = (P_1, P_2, \dots, P_n)$: vector of actor priorities.
  \item $\text{crisis\_level}(c) \in \mathbb{R}$: severity of crisis (non-crisis if $\leq -1$).
  \item $\mathrm{scale}(s) = \left( \mathrm{MAD\_score}(\mathbf{D}) \right)^{-0.25}$: Mean Absolute Deviation of avg\_pump to identify high demand actors.
  \item $ r \in [1, \infty)$
  \item $r_c \in [0, \infty)$
  \item $p_e \geq 0$
  \item $\mathbb{I}_{\{P_i > 0\}}$: indicator function (1 if $P_i > 0$, else 0)
\end{itemize}

\textbf{Quota per actor} $q_i$ is defined as:

\[
q_i =
\begin{cases}
D_i \cdot r, & \text{if } \text{crisis\_level} \leq -1 \\[12pt]
\displaystyle
\frac{A_i}
{\sum\limits_{j=1}^n A_j}
\cdot W, & \text{otherwise}

\end{cases}
\]
where allocation weights $A$ during crisis is
\[
A = D_i \cdot \exp\left(P_i \cdot c \cdot p_e \cdot s \right) \cdot \mathbb{I}_{\{P_i > 0\}}
\]
and the total available water $W$ during crisis is:
\[
W = \frac{r_c}{\max(1, c)} \cdot \sum_{k=1}^n D_k \cdot \mathbb{I}_{\{P_k > 0\}}
\]

\section{Incentives Policy Design}
This incentives policy promotes \textbf{equitable and sustainable water use} by influencing actor behavior 
through \textbf{penalties for overuse} and \textbf{rewards for conservation}. It is designed to balance three core objectives:

\begin{itemize}
    \item \textbf{Fairness:} Incentives are scaled relative to each actor’s quota and demand, ensuring proportionate responses.
    \item \textbf{Equity:} The policy accounts for actor priority, offering more leniency to high-priority users while maintaining system-wide consistency.
    \item \textbf{Ecosustainability:} Incentive strength increases with crisis severity, encouraging conservation and system resilience under stress.
\end{itemize}

\subsection{Policy Mechanism}
\subsubsection*{Core Mechanism}
The policy calculates an \textit{incentive rate} based on the difference between actual water use and the actor’s quota 
(the \textit{overuse}). This rate is scaled by:

\begin{itemize}
    \item \textbf{policy\_rate ($r_p$):} Governs the overall intensity of incentives, controlling how strongly behavior 
    is encouraged or discouraged.
    \item \textbf{equity\_spread ($eqs$):} Adjusts the baseline leniency by influencing how much the actor’s priority 
    affects the scaling of incentives. A higher equity spread means more leniency towards higher-priority actors.
    \item \textbf{crisis\_sensitivity ($cs)$:} Dictates how sharply the incentive system reacts to crisis severity. 
    When crisis levels increase, incentives become more severe, compelling actors to reduce consumption.
\end{itemize}

\subsubsection*{Subsidy and Fine Asymmetry}
The incentive function is \textbf{asymmetric} between penalties and subsidies:

\begin{itemize}
    \item The \textbf{subsidy\_fine\_ratio ($\gamma$)} controls the relative steepness of rewards (subsidies) versus penalties 
    (fines). This allows the policy to fine-tune whether conservation is rewarded more gently or harshly compared to overuse penalties.
    \item Both penalties and subsidies are bounded by maximum caps - \textbf{max\_fine\_weight ($F_{max}$)} and 
    \textbf{max\_subvention\_weight ($S_{max}$)}. This prevents excessively harsh fines or unreasonably large subsidies.
\end{itemize}

\subsubsection*{Crisis Memory and Dynamic Prioritization}
An additional important feature is the policy’s ability to \textbf{incorporate historical crisis information} to 
modulate responses dynamically:

\begin{itemize}
    \item The policy tracks whether a \textit{crisis is ongoing}, based on both current water availability and 
    consumption patterns.
    \item A parameter \(\alpha\) balances how much weight to give to \textit{past crisis conditions} versus the 
    \textit{current crisis status}.
    \item This memory effect allows the system to \textbf{anticipate future water shortages} and 
    \textbf{adjust crisis severity accordingly}. As a result, if the crisis persists over multiple iterations, the crisis level can escalate, increasing 
    the strictness of incentives.

\end{itemize}

\subsection{Mathematical Formulation}
Let:
\begin{itemize}
    \item $\mathbf{D} = (D_1, D_2, \dots, D_n)$: vector of average water pumped per actor.
    \item $\mathbf{P} = (P_1, P_2, \dots, P_n)$: vector of actor priorities.
    \item $\mathbf{C} \in \mathbb{R}$: vector of severity of crisis (non-crisis if $\leq -1$).
    \item $\mathbf{A} = (A_1, A_2, \dots, A_n)$: vector of water pumped by each actor.
    \item $\mathbf{W} = (W_1, W_2, \dots, W_t)$: vector of water flows.
    \item $\mathbf{Q} = (Q_1, Q_2, \dots, Q_n)$: vector of quotas assigned to each actor.
    \item $\mathbf{I} = (I_1, I_2, \dots, I_n)$: The average income of the actor.
    \item ${r_p} \in (0,\infty)$: The base rate of the incentive.
    \item $eqs \in (0, \infty)$: The equity spread factor.
    \item $cs \in (0, \infty)$: The crisis sensitivity factor.
    \item $\alpha in [0, 1]:$ The weight given to the current crisis versus the past crisis
    \item $\gamma \in (0, \infty)$: The ratio of subsidy to fine.
    \item $F_{max} \in (0, \infty)$: The maximum weight for fines.
    \item $S_{max} \in (0, \infty)$: The maximum weight for subsidies.
    \item \(n\): Number of actors in the system.
\end{itemize}

\noindent We define the effective $crisis\_level$ as a weighted average::
\[
crisis\_level = \alpha * next\_crisis\_level + (1 - \alpha) * C[-1]
\]
where \(C[-1]\) is the crisis level from $EcologyManager$ \\
\quad and \(\alpha \in [0,1]\) controls the weight given to the current crisis versus the past crisis.\\
\quad and  $next\_crisis\_level$ is equivalent to $EcologyManager.compute\_crisis(water\_flow, avg\_pump)$
\\ \\ 
Modify $crisis\_level$ to account for the continuing crisis:
\[
crisis\_level = \begin{cases}crisis\_level + C[-2] & \text{if }crisis\_level \geq 1 \quad or \quad C[-2] \geq 1 \\
    crisis\_level & otherwise
\end{cases}
\]
\\

\noindent We define:
\begin{itemize}
    \item $c = crisis\_level + 1$
    \item $policy\_regularization(\tau) = eqs \times \exp(-cs \times c)$.
    \item $priority\_adjustment (p_i) = \exp(2 - P_i)$
\end{itemize} 
$policy\_regularization(\tau)$ controls balance between fairness and equity, and how that balance shifts as crisis severity increases.
Penalty rates shift from relative overuse to absolute overuse as crisis level increases. 

\noindent Let the overuse rate be:
\[
    u_i = W_i - Q_i
\]
\\
\noindent The raw incentive rate before asymmetry and capping is:
\[
r_p^{i} = r_p \times \frac{u_i}{Q_i +  \tau + p_i}
\]
\\
\noindent The incentive rate is asymmetric for subsidies (underuse) and fines (overuse):
\[
r_p^{i} =
\begin{cases}
r_p^{i} \times \gamma, & \text{if } r_p^{i} < 0 \quad (\text{subsidy}) \\
r_p^{i}, & \text{if } r_p^{i} \geq 0 \quad (\text{fine})
\end{cases}
\]
\\
\noindent The incentive rate is bounded by caps to ensure fairness and predictability:
\[
r_p^{i} = \min\big(\max(r_p^{i}, S_{max}), F_{max}\big)
\]

\noindent The actual monetary incentive applied to actor $R_i$ is:
\[
R_i = I_i \times r_p^{i}
\]

\section{Policy Optimization}
Policy optimization is treated as a multi-objective optimization problem, focused on identifying optimal trade-offs 
between ecological sustainability, economic performance, and actor satisfaction. 
\\ \\ 
This process involves adjusting both quota parameters and incentive parameters to discover the best-performing policies 
within a well-defined search space. The optimization is implemented using a Pareto optimization approach within an 
evolutionary algorithm framework. The search process is powered by the off-the-shelf Optuna library, which efficiently 
explores the parameter space and constructs the Pareto front — the set of non-dominated solutions representing the best 
trade-offs among objectives.

\subsection{Implementation via PolicyRegulator Class}
The policy optimization logic is encapsulated in the $PolicyRegulator$ class, which orchestrates the following sequence:
\begin{itemize}
    \item \textbf{Define Function Arguments}: Configure the range and type of policy parameters (quota and incentive-related) to explore during optimization.
    \item \textbf{Run Fast Simulation}: Use lightweight simulations to quickly evaluate large numbers of candidate policies under varied conditions.
    \item \textbf{Objective Functions}: Simultaneously optimize ecological, economic, and satisfaction outcomes.
\[
\min \; I_{ecol}, \quad
\max \; I_{econ}, \quad
1 - S_{priority} \leq 0
\]
where: $I_{ecol}$ is the ecological impact, $I_{econ}$ is the economic impact, and $S_{priority}$ is the satisfaction of actors.
    \item \textbf{Find Pareto Front}: Identify the set of policies that offer the best possible trade-offs between the three objectives.
    \item \textbf{Select Best Policy}: Use gradient-based gain functions to rank policies within the Pareto front.
\end{itemize}

\subsection{Optimization Results}

We applied the Optuna-enabled NSGA-II evolutionary algorithm for multi-objective optimization, executing 60 trials with an initial population size of 15 while keeping all other parameters at their default settings. Convergence was observed around the 35th trial, after which improvements in the objective space plateaued.

Notably, the final run of the optimization yielded a Pareto front consisting of a single non-dominated solution. This outcome is atypical, as previous runs generally produced multiple non-dominated solutions. Despite this, the resulting solution was found to be robust and well-aligned with the objectives, clearly illustrating the trade-offs inherent in the problem formulation. The emergence of a singular dominant solution suggests that, under the given parameterization, a strong compromise exists between competing objectives, and further diversification may require broader hyperparameter exploration or relaxed constraint settings.
\begin{figure}[h]
    \centering
    \includegraphics[width=1.0\textwidth]{./assets/images/hypervolume.png}
    \caption{Hypervolume history over optimization trials. The plot illustrates the convergence behavior of the NSGA-II algorithm, with hypervolume improving steadily and stabilizing around trial 35. This indicates a progressive exploration of the objective space and eventual convergence toward a robust Pareto-optimal solution.}
    \label{fig:hypervolume_history}
\end{figure}

\subsubsection{Observations}
\begin{figure}
    \centering
    \includegraphics[width=0.7\textwidth]{./assets/images/objectives_correlation.png}
    \caption{Correlation matrix of optimization objectives. The plot highlights strong inverse relationships between key objectives, particularly between actor satisfaction and ecological impact ($r = -1.0$), and between cooperation and economic outcome ($r = -0.7$). These correlations underscore the inherent trade-offs in the multi-objective optimization landscape.}
    \label{fig:objectives_correlation}
\end{figure}

\begin{enumerate}

    \item \textbf{Satisfaction--Ecology Tradeoff}  
    
    A strong negative correlation ($r = -1.0$) was observed between actor satisfaction and ecological impact, indicating a near-perfect inverse relationship. Policies that maximize user satisfaction, particularly for high-priority actors, consistently result in elevated ecological stress, highlighting the fundamental tension between user demands and environmental preservation.

    \begin{itemize}
        \item Increasing water availability or allocation flexibility during scarcity events consistently improved satisfaction across actor types but led to higher ecological degradation.
        \item Emphasizing high-priority actor satisfaction skewed allocations, improving individual-level outcomes but intensifying overall environmental stress.
        \item Policies that constrained resource distribution in order to preserve ecological thresholds yielded the lowest satisfaction levels but were the most effective in limiting ecological impact.
        \item \textbf{Figure~\ref{fig:importance_satisfaction_ecology}} shows the hyperparameter importance plot. Priority weighting and crisis water thresholds emerge as the most influential parameters for this tradeoff, underscoring their central role in balancing satisfaction against ecological outcomes.
    \end{itemize}

    \item \textbf{Cooperation--Economy Link via Subvention}

    A moderately strong negative correlation ($r = -0.7$) was observed between economic outcomes and cooperation, particularly under subvention-heavy policy regimes. While subventions enhance individual actor incentives and promote cooperative behavior, they may also reduce systemic economic efficiency by encouraging resource redistribution over true productivity gains.

    \begin{itemize}
        \item Subventions increased cooperation by improving short-term actor returns, but they did not significantly raise total system economic output.
        \item Higher levels of cooperation under subsidy regimes often led to economic fragility and dependency, reducing long-term adaptability.
        \item Economic benefits were concentrated among actors best positioned to leverage subventions, raising concerns about equity and structural bias.
        \item \textbf{Figure~\ref{fig:importance_cooperation_economy}} illustrates that subvention-related hyperparameters are dominant drivers of cooperation. However, their marginal contribution to net economic outcomes highlights a potential misalignment in incentive structures.
    \end{itemize}

    \item \textbf{Diffusion of Cooperation Under Dissatisfaction}:  
    When satisfaction is achieved, cooperation levels tend to stabilize around a compact distribution—typically centered near 0.5—indicating a balanced and cohesive system. In contrast, low satisfaction scenarios produce more diffuse and variable cooperation, reflecting fragmentation that can destabilize system dynamics.  
    \begin{itemize}
        \item Medium-priority, high-demand actors tend to be less cooperative and more competitive during balanced cooperation-competition regimes. This arises because they occupy an intermediate position with limited buffer capacity and lower prioritization during scarcity.
        \item This dynamic is especially pronounced in the most ecologically sound solutions, where increased competition among high-demand actors plays a key role in system resilience.
    \end{itemize}

\noindent The $policy\_optimization$ notebook illustrates the difference in cooperation patterns between satisfied and unsatisfied regimes. The and $single\_scenario\_*$ demonstrate different single scenario outcomes for 
cooperation under satisfaction and dissatisfaction simulations.
\end{enumerate}

\begin{figure}[h]
    \centering
    \includegraphics[width=0.7\textwidth]{./assets/images/optimization_hyp_importances_ecol.png}
    \caption{Hyperparameter importance for the Ecological Objectives. Priority weights and water availability thresholds are key drivers of the ecological health.}
    \label{fig:importance_satisfaction_ecology}
\end{figure}

\begin{figure}[h]
    \centering
    \includegraphics[width=0.7\textwidth]{./assets/images/optimization_hyp_importances_econ.png}
    \caption{Hyperparameter importance for cooperation and economic outcomes. Subvention levels and cooperation thresholds strongly influence cooperation, while their economic effects are more diffuse.}
    \label{fig:importance_cooperation_economy}
\end{figure}

\subsection{Post-Optimization Evaluation and Manual Selection}

Following the optimization, the final run of the NSGA-II algorithm yielded a single non-dominated solution on the Pareto front. Due to the lack of multiple trade-off candidates, we manually selected and evaluated additional promising trials—specifically Trial 33, Trial 46, and a configuration derived from the best parameters of the \texttt{pol\_regulator} policy.

These parameter sets were tested using the \texttt{multi\_scenario\_analysis} notebook to assess their performance across multiple objectives. To further refine outcomes, we manually adjusted some parameters to improve simulation results, focusing on satisfaction, economic impact, and ecological impact.

Table~\ref{tab:manual_vs_optimized} presents a comparison of the evaluated parameter sets. Trial 33 demonstrated the most favorable balance, achieving absolute satisfaction while maintaining acceptable ecological and economic trade-offs. Based on this analysis, Trial 33 was selected as the final solution for submission.

\begin{table}[h]
    \centering
    \begin{tabular}{l|c|c|c|c}
        % \toprule
        \textbf{Parameter Set} & \textbf{Satisfaction} & \textbf{Violations} & \textbf{Economic Impact} & \textbf{Ecological Impact} \\ \hline
        % \midrule
        Trial 33 (\textit{best})           & \textbf{1.000 }  & 0              & \textbf{0.885}                   & \textbf{0.935 }                    \\
        Trial 46                           & \textbf{1.000}   & 0              & 0.840                             & 0.941                     \\
        % \bottomrule
    \end{tabular}
    \caption{Comparison of parameter sets based on key optimization objectives}
    \label{tab:manual_vs_optimized}
\end{table}

\section{Conclusion}
The optimization revealed fundamental trade-offs between actor satisfaction, ecological impact, cooperation, and economic outcomes. Notably, a perfect negative correlation between satisfaction and ecological health underscores the challenge of balancing human needs with environmental sustainability. Cooperation improves with higher subvention levels but exhibits a negative correlation with economic performance, suggesting diminishing returns and potential long-term fragility when over-reliant on subsidies.
\\ \\
A key behavioral insight is the pattern of cooperation relative to satisfaction: cooperation tends to concentrate around moderate levels (approximately 0.5) when satisfaction is high, indicating a stable and balanced system. In contrast, dissatisfaction leads to more diffuse and fragmented cooperation patterns, reflecting instability and less cohesive interactions.
\\ \\
These insights are further validated through detailed single scenario analyses, reinforcing the robustness of the observed trade-offs and behavioral dynamics. Overall, the study highlights the necessity for nuanced policy designs that carefully balance competing objectives and leverage evolutionary optimization to identify both effective solutions and meaningful system behaviors.
\end{document}